% Generated draft: requires mathematical and source review.
\documentclass[11pt]{article}
\usepackage[margin=25mm]{geometry}
\usepackage{fontspec}
\setmainfont{Latin Modern Roman}
\usepackage{amsmath,amssymb,mathtools,unicode-math}
\setmathfont{Latin Modern Math}
\usepackage{enumitem,needspace,xurl,hyperref}
\setlength{\parindent}{0pt}
\setlength{\parskip}{5pt}
\setlength{\emergencystretch}{3em}
\newcommand{\N}{\mathbb{N}}
\newcommand{\Z}{\mathbb{Z}}
\newcommand{\Q}{\mathbb{Q}}
\newcommand{\R}{\mathbb{R}}
\newcommand{\C}{\mathbb{C}}
\newcommand{\F}{\mathbb{F}}
\newcommand{\A}{\mathbb{A}}
\newcommand{\PP}{\mathbb{P}}
\newcommand{\E}{\mathbb{E}}
\newcommand{\eps}{\varepsilon}
\newcommand{\dd}{\,\mathrm{d}}
\begin{document}
\section*{UnsolvedMath ID: 30000271}
\textbf{Source problem code:} catalogue entry
\subsection{Definitions and mathematical statement}
The exact extracted request has two parts: find further examples in tire-track geometry, and determine the two-dimensional analogue of the ``uniform cylindrical logarithm''. The latter phrase is retained exactly but is not defined; despite the verification field calling the sentence coherent, it does not identify a mathematical object. The supplied assessment connects this part to a two-dimensional flotation problem. No correction of ``logarithm'' or substitution of a particular flotation conjecture is made silently.

For orientation, a planar bicycle configuration with wheelbase $\ell>0$ can be represented by rear and front paths $r,f$ satisfying
\[
 f(t)-r(t)=\ell v(t),\qquad |v(t)|=1,\qquad
 r'(t)\parallel v(t),
\]
where the last condition describes the rear wheel's no-side-slip constraint. The relevant closed-curve ambiguity condition, regularity, equivalence of tracks and examples already known are absent from the excerpt. Thus ``further examples'' cannot be turned into an all-examples classification or a unique algebraic condition without its source antecedent.

A possible planar flotation interpretation would concern a homogeneous body $K\subset\R^2$ of density ratio $0<\rho<1$. For each upward unit vector $n$, a waterline $x\cdot n=h_n$ cuts off a submerged part $K_n$ of area $\rho|K|$; torque balance requires the centroids of $K_n$ and $K$ to be aligned with $n$. Demanding this for every orientation is an all-orientations equilibrium condition, distinct from stability of each equilibrium. The source does not supply $K$, its convexity/smoothness class, density, orientation quantifier or a definition connecting that condition to its ``cylindrical logarithm'' phrase. This is only an explicit contextual possibility suggested by the assessment.

The supplied literature triage checked 21 August 2026 reports additional ambiguous closed bicycle curves and integrable examples, but no full classification or resolution of the associated planar flotation question. Those partial results do not fix the missing object or require the original examples request to be a universal classification theorem.
\subsection{Short English statement}
The entry seeks more tire-track examples and an undefined planar analogue of a ``uniform cylindrical logarithm''. The supplied assessment points to bicycle-curve ambiguity and flotation, but the exact geometric conditions and terminology are missing. Later examples are reported as partial progress.
\subsection{Sources}
\begin{itemize}
\item Local UnsolvedMath registry, permanent ID 30000271, exact source code catalogue entry; \emph{Tire-Track Geometry and Cylindrical Logarithms}. Proposed year 2005, status partially\_solved, exact wording including ``uniform cylindrical logarithm'' retained.
\item Supplied report citation: \emph{Dynamical Systems (2006)}, Oberwolfach Reports, DOI 10.4171/owr/2005/31, pp. 1743--1798, \url{https://doi.org/10.4171/owr/2005/31}. Title year differs from DOI/proposed year.
\item Supplied literature triage checked 2026-08-21: ``original Oberwolfach report'', \url{https://doi.org/10.4171/owr/2005/31}; ``Tire tracks and integrable curve evolution'', \url{https://arxiv.org/abs/1705.06314}. This dated OWR-derived curation requires independent verification of time-sensitive status and cited claims.
\end{itemize}
\end{document}
